\Article{17}
The degree of confidence to be placed in this determination must be judged by the mean error to be feared either in the determination itself or in its square; in the latter case, this mean error is the square root
\SourceLeaf{081}
of the mean value of the expression
\[
\left(\frac S\sigma-\mu\mu\right)^2
\]
whose expansion follows by reasoning like that set out in the preceding paper, Arts. 39 ff. For the sake of brevity, that reasoning is omitted here, and we give only the formula itself. Namely, the mean error to be feared in the determination of the square \(\mu\mu\) is expressed by
\[
\BareRoot{\left(\frac{2\mu^4}{\sigma}+\frac{\nu^4-3\mu^4}{\sigma\sigma}\cdot N\right)}
\]
where \(\nu^4\) denotes the mean value of the fourth powers of errors. whose weight is \(=1.\), and \(N\) the sum
\[
\begin{aligned}
&(a\alpha+b\beta+c\gamma+\text{etc.})^2+(a'\alpha'+b'\beta'+c'\gamma'+\text{etc.})^2\\
&\qquad+(a''\alpha''+b''\beta''+c''\gamma''+\text{etc.})^2\text{etc.}
\end{aligned}
\]
In general this sum cannot be reduced to a simpler form; but, as can be shown in the same way as in Art. 40 of the preceding paper, its value always lies between \(\pi\) and \(\frac{\sigma\sigma}{\pi}\). Under the hypothesis on which the theory of least squares was founded from the outset, the term containing this sum vanishes, since \(\nu^4=3\mu^4\); and the precision to be attributed to the mean error determined by the formula \(\BareRoot{\frac S\sigma}\) is the same as if that mean error had been derived from \(\sigma\) exactly known errors according to Arts. 15 and 16 of the preceding paper.

\Article{18}
As we have seen above, two things are required for the adjustment of observations: first, the correlates of the conditional equations, that is, the numbers \(A,B,C\) etc. satisfying equations (12), must be derived; second, these numbers must be substituted in equations (10). An adjustment obtained in this way may be called \textit{perfect} or \textit{complete}, to distinguish it from an \textit{imperfect} or \textit{incomplete} adjustment. By the latter term we shall designate adjustments that result from the same equations (10), but with values of the quantities \(A,B,C\) etc. substituted that do not satisfy equations (12), that is, values satisfying only some of them or none. As for changes to the observations that cannot be expressed by formulas (10), they are excluded both from the present investigation and from the term “adjustments.” Since, insofar as equations (10) apply, equations (13) are wholly equivalent to equations (12), the distinction may also be stated as follows: \SourceLeaf{082}
Completely adjusted observations satisfy all the conditional equations \(X=0,\allowbreak Y=0,\allowbreak Z=0\) etc.; incompletely adjusted observations satisfy either none or at least not all of them. An adjustment, therefore, that satisfies all the conditional equations is necessarily complete.

\Article{19}
Since it follows from the very notion of an adjustment that the combination of two adjustments again constitutes an adjustment, it is readily seen that it makes no difference whether the rules by which a perfect adjustment is obtained are applied directly to the original observations or to observations that have already been incompletely adjusted.

Indeed, let \(-\theta,-\theta',-\theta''\) etc. be a system of incomplete adjustments obtained from formulas (I):
\[
\begin{aligned}
\theta p&=A^0a+B^0b+C^0c+\text{etc.}\\
\theta'p'&=A^0a'+B^0b'+C^0c'+\text{etc.}\\
\theta''p''&=A^0a''+B^0b''+C^0c''+\text{etc.}\\
&\text{etc.}
\end{aligned}
\]
Since the observations altered by these adjustments are assumed not to satisfy all the conditional equations, let \(\mathfrak A^*,\mathfrak B^*,\mathfrak C^*\) etc. be the values that \(X,Y,Z\) etc. take upon their substitution. We must find numbers \(A^*,B^*,C^*\) etc. satisfying equations (II):
\[
\begin{aligned}
\mathfrak A^*&=A^*[aa]+B^*[ab]+C^*[ac]+\text{etc.}\\
\mathfrak B^*&=A^*[ab]+B^*[bb]+C^*[bc]+\text{etc.}\\
\mathfrak C^*&=A^*[ac]+B^*[bc]+C^*[cc]+\text{etc.}
\end{aligned}
\]
etc. Once this has been done, a complete\textperiodcentered\ adjustment of the observations thus altered is effected by the new changes \(-\kappa,-\kappa',-\kappa''\) etc., where \(\kappa,\kappa',\kappa''\) etc. are computed from formulas (III):
\[
\begin{aligned}
\kappa p&=A^*a+B^*b+C^*c+\text{etc.}\\
\kappa'p'&=A^*a'+B^*b'+C^*c'+\text{etc.}\\
\kappa''p''&=A^*a''+B^*b''+C^*c''+\text{etc.}
\end{aligned}
\]
Let us now inquire how these corrections are related to the complete adjustment of the original\SourceLeaf{083}
observations. First, it is evident that
\[
\begin{aligned}
\mathfrak A^*&=\mathfrak A-a\theta-a'\theta'-a''\theta''-\text{etc.}\\
\mathfrak B^*&=\mathfrak B-b\theta-b'\theta'-b''\theta''-\text{etc.}\\
\mathfrak C^*&=\mathfrak C-c\theta-c'\theta'-c''\theta''-\text{etc.}
\end{aligned}
\]
etc. Substituting in these equations the values of \(\theta,\theta',\theta''\) etc. from (I), as well as the values of \(\mathfrak A^*,\mathfrak B^*,\mathfrak C^*\) etc. from (II), we find
\[
\begin{aligned}
\mathfrak A&=(A^0+A^*)[aa]+(B^0+B^*)[ab]+(C^0+C^*)[ac]+\text{etc.}\\
\mathfrak B&=(A^0+A^*)[ab]+(B^0+B^*)[bb]+(C^0+C^*)[bc]+\text{etc.}\\
\mathfrak C&=(A^0+A^*)[ac]+(B^0+B^*)[bc]+(C^0+C^*)[cc]+\text{etc.}
\end{aligned}
\]
etc.; hence it follows that the correlates of the conditional equations satisfying equations (12) are
\[
A=A^0+A^*,\quad B=B^0+B^*,\quad C=C^0+C^*\ \text{etc.}
\]
Equations (10), (I), and (III) therefore show that
\[
\varepsilon=\theta+\kappa,\quad\varepsilon'=\theta'+\kappa'\quad\varepsilon''=\theta''+\kappa''\ \text{etc.}
\]
that is, the complete adjustment of the observations is the same whether computed directly or indirectly, starting from an incomplete adjustment.